\newcommand{\hmwkTitle}{Homework 3}
\newcommand{\hmwkDueDate}{Saturday, January 30, 2016}
\newcommand{\hmwkHeaderTitle}{ICCS206 Discrete Mathematic: Homework 3}
\documentclass{article}

\usepackage{fancyhdr}
\usepackage{amsmath}
\allowdisplaybreaks
\usepackage{amsthm}
\usepackage{amsfonts}
\usepackage{tikz}
\usepackage[plain]{algorithm}
\usepackage{algpseudocode}
\usepackage{enumitem}
\usepackage{graphicx}
\usepackage{hyperref}

\usetikzlibrary{automata,positioning}

%
% Basic Document Settings
%

\topmargin=-0.45in
\evensidemargin=0in
\oddsidemargin=0in
\textwidth=6.5in
\textheight=9.0in
\headsep=0.25in

\linespread{1.1}

\pagestyle{fancy}
\lhead{\hmwkAuthorName}
\chead{} % empty
\rhead{\hmwkHeaderTitle}
\cfoot{\thepage}

\renewcommand\headrulewidth{0.4pt}
\renewcommand\footrulewidth{0pt}

\setlength\parindent{0pt}

%
% Homework Details
%

\providecommand{\hmwkTitle}{Homework}
\providecommand{\hmwkDueDate}{TBD}
\providecommand{\hmwkClass}{ICCS206 Discrete Mathematic}
\providecommand{\hmwkClassTime}{Section\ 1}
\providecommand{\hmwkClassInstructor}{Pakawut Jiradilok}
\providecommand{\hmwkAuthorName}{\textbf{Jiraroj Wiruchpongsanon (6781617)}}
\providecommand{\hmwkHeaderTitle}{\hmwkClass:\ \hmwkTitle}

%
% Title Page
%

\title{
    \vspace{2in}
    \textmd{\textbf{\hmwkClass:\ \hmwkTitle}}\\
    \normalsize\vspace{0.1in}\small{Due\ on\ \hmwkDueDate}\\
    \vspace{0.1in}\large{\textit{\hmwkClassInstructor\ \hmwkClassTime}}
    \vspace{3in}
}

\author{\hmwkAuthorName}
\date{}

\renewcommand{\part}[1]{\textbf{\large Part \Alph{partCounter}}\stepcounter{partCounter}\\}

%
% Helper Commands
%

\newcounter{partCounter}
\newcommand{\solution}{\textbf{\large Solution}}

\newtheorem*{theorem}{Theorem}

\theoremstyle{remark}
\newtheorem*{remark}{Remark}

\begin{document}

\maketitle
\newpage


\section*{Problem 1}
A chocolate-breaking game:
\begin{enumerate}[label=\arabic*.]
    \item Your friend picks the size of a chocolate bar (e.g., $5 \times 2$ or $11 \times 11$).
    \item After seeing the size, you choose whether you or your friend makes the first break.
    \item Players alternate breaking one piece at a time until all pieces are $1 \times 1$.
    \item Whoever breaks the last piece wins.
\end{enumerate}
Find and justify the winning strategy. (Bonus: try it with friends.)

\medskip
\solution

we try to solve this using hard induction, splitting this two dimentional
problem into quantitative one dimentional one, instead of trying to prove
the claim of

\[
\text{it takes } mn - 1 \text{ breaks until all pieces of } m \times n
\text{ chocolate bar are } 1 \times 1 \text{ pieces}
\]

this hard, so we rewrite the claim as

\[
P(k) : \text{it takes } k - 1 \text{ breaks until all pieces of } k
\text{ pieces chocolate are } 1 \times 1 \text{ pieces}
\]

for the base case where $n_0 = 1$

\[
\begin{aligned}
P(1) &= 0 \\
P(2) &= 1 \\
P(4) &= 3 && \text{in case of } 2 \times 2 \text{ bar}
\end{aligned}
\]

then, we assume P(k) for arbitrary $(k \ge n_0)$, it'll take k - 1 breaks.
we do one break anywhere, the bar will split up into a smaller bar (than k)
of q and p --- p + q = k.

\vspace{1em}

from our inductive hypothesis, we're given that a bar of size p, takes p - 1
breaks until it's all 1 * 1 pieces, and a bar of size q, takes q - 1 breaks
until it's all 1 * 1 pieces. so ...

\[
\begin{aligned}
\text{breaks} &= 1 + (p - 1) + (q - 1) && \text{\{first break from k\} +
\{inductive of it's result\}} \\
&= p + q - 1 && k = p + q \\
&= k - 1
\end{aligned}
\]

the inductive case is satisfied, therefore, we'll need k - 1 breaks to
finish (achieve all 1 * 1 pieces) a k sized bar of chocolate.

\vspace{1em}

our strategy could be to go first if there were even pieces, and go second
if there were odd pieces --- for a game of two players.

\vspace{1em}
\hrulefill

\newpage

\section*{Problem 2}
Consider an $n$-piece jigsaw puzzle assembled by repeatedly joining two existing blocks (starting from single pieces). Prove that assembling the puzzle takes exactly $n-1$ steps.

\medskip
\solution

\begin{proof}

we claim that

\[
P(n) : \text{it takes } n - 1 \text{ steps to complete an } n\text{-piece
jigsaw}
\]

for the \textbf{base case} where $n_0 = 1$

\[
\begin{aligned}
P(1) &= 0 && \text{we didn't join any pieces yet} \\
P(2) &= 1 \\
P(3) &= 2
\end{aligned}
\]

by observing this series, we got

\[
P(k + 1) = P(k) + 1 = k
\]

since we merge a piece of jigsaw in every iteration, and the jigsaw will be
complete when $P(n) = n - 1$

\[
P(k) = P(k - 1) + 1 = k - 1
\]

so

\[
\begin{aligned}
P(k + 1) &= 1 + P(k) \\
&= 1 + (k - 1) \\
&= k
\end{aligned}
\]

we got $P(k + 1) = k$, which could be formatted to $\boxed{P(n) = n - 1}$.
    
\end{proof}

\vspace{1em}
\hrulefill

\section*{Problem 3}
Suppose Brew\&Bev gift certificates come in 7-Baht and 4-Baht values, with unlimited supply. Show that any price $p \ge 18$ can be paid with exact change. (Hint: check 22 and 25 Baht.)

\medskip
\solution

\begin{proof}

we split this proof into four congruence cases, prove each of them by
observation, we claim that

\[
\text{any integer } p \text{ that's at least } 18 \text{ can be constructed
with the sum of } 7 \text{ and } 4
\]

classify the integers by their residue modulo 4, given the base integer
of each category (residue of 0, 1, 2 and 3), sum it up with arbitrary $4n$
where $n$ is also an integer --- we could construct any integer following
that base integers.

\vspace{1em}

so, now we just need to find those 4 base integers, intersect their range
--- which is $b - 4$ where $b$ is that base integer.

\begin{enumerate}
\item case 1: $b_0 \equiv 0 \pmod 4$

\vspace{1em}

we got the smallest $b_0$ as $0$, which we can now construct any integer of
this same residue

\[
0 = 0 \cdot 7 + 0 \cdot 4
\]

\[
0 \equiv 0 + 4n \pmod 4
\]

and thus, the range of this base residue case is
\[
(-4, \infty) \text{ (around } (0, \infty)).
\]

\item case 2: $b_1 \equiv 1 \pmod 4$

\vspace{1em}

the smallest $b_1$ following the construction rule is
\[
21 = 7 + 7 + 7,
\]
hence every integer of this residue is

\[
21 \equiv 21 + 4n \pmod 4, \qquad 21 + 4n = 7 + 7 + 7 + 4n
\]

and the range is
\[
(21 - 4, \infty) = (17, \infty).
\]

\item case 3: $b_2 \equiv 2 \pmod 4$

\vspace{1em}

the smallest $b_2$ is
\[
14 = 7 + 7,
\]
so any integer of this residue satisfies

\[
14 \equiv 14 + 4n \pmod 4, \qquad 14 + 4n = 7 + 7 + 4n
\]

and the range is
\[
(14 - 4, \infty) = (10, \infty).
\]

\item case 4: $b_3 \equiv 3 \pmod 4$

\vspace{1em}

the smallest $b_3$ is
\[
7 = 7,
\]
giving

\[
7 \equiv 7 + 4n \pmod 4, \qquad 7 + 4n = 7 + 4n
\]

and the range is
\[
(7 - 4, \infty) = (3, \infty).
\]
\end{enumerate}

we then intersect the range from those cases

\[
(0, \infty) \cap (17, \infty) \cap (10, \infty) \cap (3, \infty) = (17,
\infty) \approx [18, \infty)
\]

thus $p$ is true for all integers over $18$.\vspace{1em}

\end{proof}

\vspace{1em}
\hrulefill

\section*{Problem 4}
Show that every positive integer is a product of a power of two and an odd integer.

\medskip
\solution

we claim that

\[
P(x) : \forall x \in I^+ \mid x = 2^n \cdot m \text{ where } n \ge 0
\text{ and } m \in \mathbb{O}
\]

for the \textbf{base case} where $n_0 = 1$

\[
\begin{aligned}
P(1) &= 2^0 \cdot 1 \\
P(2) &= 2^1 \cdot 1 \\
P(3) &= 2^0 \cdot 3
\end{aligned}
\]

then, we assume P(k) for arbitrary $(k \ge n_0)$, we prove $P(k+1) = 2^{n'}
\cdot m'$ where $n \ge 0$ and $m \in \mathbb{O}$ --- there'll be two case
here

\begin{enumerate}
\item case 1: $k + 1 \in \mathbb{O}$

for this case

\[
k + 1 = 2^0 \cdot (k + 1)
\]

since $(k + 1)$ is already an odd number, the induction hypothesis is
satisfied

\vspace{1em}

\item case 2: $k + 1 \in \mathbb{E}$

since $k + 1$ is an even integer, it could be write in a form of $k + 1 =
2k'$, and since by the induction hypothesis

\[
k' = 2^{n'} \cdot m'
\]

so

\[
\begin{aligned}
k + 1 &= 2(2^{n'} \cdot m') \\
&= 2^{n' + 1} \cdot m'
\end{aligned}
\]
\end{enumerate}

thus, P(k+1) holds, the \textbf{induction case} is satisfy, and we can
conslude that the given statement is true.

\vspace{1em}
\hrulefill

\section*{Problem 5}
Let $T_1 = T_2 = T_3 = 1$ and 

\[T_n = T_{n-1} + T_{n-2} + T_{n-3} \text{ for } n \ge 4\] 

Show that $T_n < 2^n$.

\medskip
\solution
\begin{proof}
    
we claim that

\[
P(n) : \forall n \ge 4, ; T_n < 2^n
\]

for the \textbf{base case} where $n_0 = 4$

\[
\begin{aligned}
P(4) &= 1 + 1 + 1 = 3 < 2^4 = 16 \\
P(5) &= T_5 = T_4 + T_3 + T_2 = 3 + 1 + 1 = 5 < 2^5 = 32 \\
P(6) &= T_6 = T_5 + T_4 + T_3 = 5 + 3 + 1 = 9 < 2^6 = 64
\end{aligned}
\]

then, we assume P(k) for arbitrary $k \ge n_0$, since

\[
T_k = T_{k-1} + T_{k-2} + T_{k-3}
\]

and

\[
\begin{aligned}
T_{k-1} &< 2^{k-1} \\
T_{k-2} &< 2^{k-2} \\
T_{k-3} &< 2^{k-3}
\end{aligned}
\]

thus

\[
\begin{aligned}
T_k &< 2^{k-1} + 2^{k-2} + 2^{k-3} \\
&< 2^{k-3}(2^2 + 2^1 + 2^0)
\end{aligned}
\]

while

\[
2^k = 2^3 (2^{k-3})
\]

and since

\[
2^3 (2^{k-3}) > 7 (2^{k-3})
\]

we can conclude that the statement: $2^n > T_n$ is true.
\end{proof}
\vspace{1em}
\hrulefill

\section*{Problem 6}
Show that every integer $n \ge 1$ can be written as a sum of distinct powers of two.

\medskip
\solution

\begin{proof}

we claim that

\[
P(n) : \forall n \in I^+ \text{ can be written as a sum of distinct powers
of two}
\]

for our \textbf{base case} where $n_0 = 1$

\[
\begin{aligned}
P(1) &= 2^0 \\
P(2) &= 2^1 \\
P(3) &= 2^1 + 2^0
\end{aligned}
\]

even though it's quite obvious that this is just base 2 number system, which
you could just construct $P(k+1)$ by just adding $1_2$ to $P(k)$, I'll not
do that. we assume $P(k)$ for arbitrary $(k \ge n_0)$, let $C_k$ notates the
combination of those distinct powers of two that summed up to $k$, and try
to prove $P(k+1)$ by case as

\begin{enumerate}
\item case 1: $k + 1 \in \mathbb{O}$

for this one, since $k + 1$ is an odd integer, $k$ would be even, which
means that $C_k$ set doesn't contain $2^0$, and thus

\[
C_{k + 1} = C_k + 2^0
\]

\item case 2: $k + 1 \in \mathbb{E}$

for this one, since $k + 1$ is an even integer, it means that $k + 1$ can be
written in a form of

\[
k + 1 = 2k'
\]

and from our inductive hypothesis, $C_{k'}$ should already satisfied
$P(k')$, therefore $C_{k+1}$ could be construct by multiplying $C_{k'}$ by
$2$ (or bitwise shift to the left in binary sense) as

\[
C_{k+1} = 2 \cdot C_{k'} = {2c \mid c \in C_{k'}}
\]

which the $2^0$ place is guarantees non-existance.
\end{enumerate}

since, $P(k+1)$ is unique and representable, the \textbf{induction case}
holds, and we could conclude that the claim is true

\end{proof}

\vspace{1em}
\hrulefill

\newpage 

\section*{Problem 7}
Consider an infinite grid. Start with one unit square whose perimeter is 4 (even). Each subsequent square must share at least one edge with the existing cluster. Prove that the perimeter length remains even regardless of placement order.

\medskip
\solution

\begin{proof}
    
we claim that

\[
P(n) : \text{the perimeter of any connected figure of } n \text{ unit
squares is even}
\]

for the \textbf{base case} where $n_0 = 1$

\[
P(1) = 4
\]

then, we assume $P(k)$ for arbitrary $(k \ge n_0)$, the perimeter of
every connected $k$-square figure is even. 

\vspace{1em}
to build $P(k+1)$, add one more
unit square; it contributes four fresh edges, but since it must share at
least one edge and at most four edges with the existing figure, the net
change is

\[
P(k+1) = P(k) + 4 - 2n, \qquad \text{ when } 1 \le n \le 4
\]

because $P(k)$ is even, let $P(k) = 2q$; then

\[
P(k+1) = 2q + 4 - 2n = 2(q + 2 - n)
\]

which is still even. thus, $P(k+1)$ holds, the \textbf{induction case} is
satisfied, and we conclude the claim is true.
\end{proof}

\vspace{1em}
\hrulefill

\section*{Problem 8}
Is the following proposition true? If not, explain the flaw in the proof.

\vspace{1em}

Theorem: If we have $n$ straight line in a 2 dimensional plane $(n \ge
2)$, and none of the line are parallel to each other, then all lines must
intersect at exactly one and the same point. 

\vspace{1em}

Proof: by induction-ish?.

\vspace{1em}

Inductive Predicate:$P (n) = $ every set of $n$ straight lines intersect
at exactly one point. 

\vspace{1em}

Base Case: $n = 2$. Every two lines intersect at one
point. Therefore, $P (2)$ is true. 

\vspace{1em}

Inductive Step: Assume that every $n$ straight lines intersect at one point, we want to show that every $n + 1$
line intersect at one point.

\vspace{1em}
\begin{itemize}
\item Let the set of $n + 1$ be ${a_1, a_2, a_3, \ldots , a_n, a_{n+1}}$
\item From the inductive hypothesis ${a_1, a_2, a_3, \ldots , a_n}$ must
intersect at one point since it is a set of $n$ lines. Let us call the
intersection point for these lines $D$.
\item Similarly, ${a_2, a_3, a_4 \ldots , a_n, a_{n+1}}$ intersect at
exactly one point since there are $n$ lines in this set. Let us call this
point $E$.
\item Since both $D$ and $E$ are the point where $a_2$ intersects $a_3$. $D$
and $E$ are therefore the same point.
\item Therefore, ${a_1, a_2, a_3, \ldots , a_n, a_{n+1}}$ intersect at
exactly one point.
\end{itemize}

\medskip
\solution

The two set of size $n$ derived from the inductive hypothesis: $(a_1, a_2,
\ldots, a_n)$ and $(a_2, a_3, \ldots, a_{n+1})$ still doesn't cover the
intersection of $a_1$ and $a_{n+1}$, so we have to do a quick prove on these
two lines as

\vspace{1em}

\begin{proof}
since line $a_1$ and $a_{n+1}$ are not parallel, they will have to intersect
at some point in a 2d plane
\end{proof}

and this completes the induction case of our prove, and we can confidently
conclude that our inductive predicate holds true
\end{document}
